\scanpage{139}
% archive sheet number: 138

Choisissons bases $\unsure{1(,u)|(1,u)}$ de $H$, $(1,v)$ de $H'$
(d'où $b,c,b',c'$), $\varphi$ et $\varphi'$ sont données par

\[
\varphi(u)=z,\qquad\varphi'(v)=z'
\]

avec $z\in k$

\[
z^2-bz+c=0\qquad z'{}^2-b'z+c'=0,
\]

et on cherche $t$ solution de

\[
t^2-bb't+(b^2b'+b'{}^2b-4cc')=0.
\]

Mais le discriminant de cette équation est

\[
\Delta=(\deleted{\unsure{b}}bb')^2-4\bigl(b^2\deleted{b}c'+b'{}^2\deleted{b}c-4cc'\bigr)=
\underbrace{(b^2-4c)}_{\delta}\underbrace{(b'{}^2-4c')}_{\delta'}
\]

Or \deleted{\unsure{comme}} (\unsure{se|en|les} plaçant \deleted{en car. 2} heuristiquement
d'abord sur un corps de car $\ne2$)

\[
2z=b\pm\sqrt{\delta},\qquad 2z'=b'\pm\sqrt{\delta'}
\]

\[
\text{\deleted{donc}}\qquad\delta=(2z-b)^2\qquad\delta'=(2z'-b)^2
\]

(vrai sans condition corps ou caract.)
et les solutions de l'équation en $t$ sont \inserted{données} par
(sur corps car $\ne2$) \unsure{d}

\[
\deleted{\unsure{\tfrac14}}\,(2t-bb')^2=\Delta=\delta\delta'=\bigl((2z-b)(2z'-b')\bigr)^{\unsure{2}}
\]

donc on peut \deleted{poser} supposer

\[
2t-bb'=\pm(2z-b)(2z'-b')
\]

soit

\[
2t=bb'\pm(2z-b)(2z'-b').
\]

Choisissant $-$ des \deleted{\unsure{sig}} les $\pm$, on trouve

\[
2t=\deleted{\unsure{4}}-4zz'+2(\deleted{\unsure{4}}b'z+bz')\deleted{\unsure{+}}\qquad\text{\deleted{soit}}
\]

\[
\boxed{t=(b'z+bz')-2zz'}
\]

l'autre racine sera

\[
t^{\unsure{1|'}}=\bigl(bb'-(b'z+bz')\bigr)+2zz'\ (=bb'-\deleted{\unsure{z}}t)
\]

